\subsection{Three different coordinate quantities}

The single-site conditional-variance influence is
\[
 H_i^\mu(f)=\mathbb E_\mu\operatorname{Var}(f\mid X_{-i})
 =\mathbb E_\mu\left[\operatorname{Var}(X_i\mid X_{-i})(D_if)^2\right].
\]
For Curie--Weiss at finite $\beta\ge0$,
\[
 \operatorname{sech}^2(\beta)I_i^\mu(f)
 \le H_i^\mu(f)\le I_i^\mu(f),
\]
because $\mathbb E[X_i\mid X_{-i}]=\tanh(\beta\sum_{j\ne i}X_j/n)$. Both definitions therefore have the same fixed-temperature threshold.

The variance explained by \emph{revealing} a coordinate,
\[
 R_i^\mu(f)=\operatorname{Var}_\mu\bigl(\mathbb E_\mu[f\mid X_i]\bigr),
\]
is different again. For zero-field Curie--Weiss and $M=n^{-1}\sum_iX_i$, exchangeability and sign symmetry give
\[
 \mathbb E[M\mid X_i]=\operatorname{Var}(M)X_i,
 \qquad R_i^\mu(M)=\operatorname{Var}(M)^2,
 \qquad I_i^\mu(M)=n^{-2}.
\]
Above the transition the reveal quantity stays bounded away from zero while flip influence vanishes. Sampling $O(\epsilon^{-2})$ coordinates estimates $M$ to additive error $\epsilon$ with constant success probability on every fixed input. Thus magnetization disproves an influence inequality without demonstrating quantum advantage.
